% FUSAO (26/10/2026): float copied unchanged from tables_C*.tex, tables_main.tex or tables_appendix.tex.
\begin{table}[htbp]\centering
\caption{Change between 2014 and 2017 in the waiting-list stock and in flow measures, by the month in which the count or window ends}
\label{tab:stab}
% Gerado por 02_audit/build20_tables.py (26/10/2026). Nao editar a mao.
\footnotesize
\fitbox{%
\begin{tabular}{lrrrrr}
\toprule
Measure, change 2014 to 2017 & March & June & September & December & Dec./Sep. \\
\midrule
Registered demand, snapshot at month end (\%) & $-$20.5 & $-$18.9 & $-$17.5 & $-$53.2 & 3.03 \\
New first requests, 12 months to month end (\%) & $-$8.3 & $-$5.4 & $-$8.2 & $-$6.2 & 0.75 \\
Median age at first request, same window (months) & $-$5.2 & $-$4.6 & $-$4.1 & $-$3.8 & 0.93 \\
Referred within 365 days, same window (share) & $+$0.130 & $+$0.146 & $+$0.172 & $+$0.191 & 1.11 \\
\bottomrule
\end{tabular}}

\tabnote{\emph{Sources:} SME publications (registered demand); SME childcare request record, first requests from July 2012, with referral dates through November 2019 (flow measures). \emph{Unit:} S\~ao Paulo municipality. Registered demand is the count in the publication for the stated month (December extraction dates in Table~\ref{tab:dates}; the September 2017 publication prints no extraction date). New first requests, median age at first request and the share of first requests referred within 365 days are measured over the twelve months ending in the stated month. Entries are changes between the 2014 and 2017 values, in per cent for registered demand and new requests, in months for median age, and as a change in the share for referral. Dec./Sep.\ denotes the change measured in December divided by the change measured in September. The comparison starts in 2014 because March 2014 is the first March whose preceding twelve-month window lies wholly inside the window of complete coverage; over 2013--2017, the stock ratio is 3.47. The referral share is the simple share of all first requests in each window referred within 365 days, without the competing-event treatment, so its level differs from that in Table~\ref{tab:service}. Part of the contrast between the stock and the flows is mechanical: the stock is a count on a single day, whereas the flow measures cover twelve months. The mean of the four published counts in the twelve months ending in the stated month changes by $-20.8$, $-20.8$, $-21.5$ and $-25.4$ per cent (Dec./Sep.\ 1.18), so most of the month dependence of the stock is a property of the single-day count (Section~\ref{sec:indicators}). The comparisons are descriptive.}
\end{table}
